\uuid{NIQ7}
\titre{Points stationnaires}
\auteur{hanani ; mzari}
\organisation{exo7}
\licence{CC BY-NC-SA 4.0}
\datecreate{2020-01-20}
\niveau{L1}
\module{Géométrie}
\chapitre{Courbes planes}
\sousChapitre{Courbes paramétrées}
\difficulte{3}
\theme{Courbes paramétrées, points stationnaires}
\exo7qcmid{lille-2.courbes.008}

\contenu{
\begin{question}
\qclassification{244.01}
\qtags{motcle=points stationnaires, tangentes}
La trajectoire $\Gamma$ d'une particule en mouvement est donnée par les équations
$$x=1+\cos t\quad \mbox{et} \quad y=t-\sin t.$$
\begin{answers}
\bad{$\Gamma$ est symétrique par rapport à l'axe des $y$.}
\bad{$\Gamma$ admet un point double.}
\good{Le point de paramètre $t=0$ est un point de rebroussement de première espèce.}
\good{La tangente à $\Gamma$ au point de paramètre $t=0$ est horizontale.}
\end{answers}
\begin{explanations}
D'abord, $x$ est paire et $y$ est impaire, donc $\Gamma$ est symétrique par rapport à l'axe des $x$. Ensuite, en notant $f(t)=(x(t),y(t))$, la résolution du système $f(t_1)=f(t_2)$ donne $t_1=t_2$, donc $\Gamma$ ne possède pas de point double. Enfin, on a :
$$f(t)=(2,0)+t^2(-1/2,0)+t^3(0,1/6)+t^3\left(\varepsilon_1(t),\varepsilon_2(t)\right)\mbox{ avec }\lim _{t\to 0}\varepsilon_i(t)=0.$$
Donc le point de paramètre $t=0$ est un point de rebroussement de première espèce et la tangente en ce point est horizontale.
\end{explanations}
\end{question}
}
